\subsection{Justificación del uso de aprendizaje profundo en el proyecto}

A partir de la comparación entre el aprendizaje automático tradicional y el aprendizaje profundo, resulta posible identificar con mayor claridad por qué este último constituye el enfoque más adecuado para el problema planteado en el presente trabajo.

En tareas de clasificación de imágenes, la información de entrada no se encuentra organizada como un conjunto reducido de variables explícitas, sino distribuida en una estructura visual compleja que incluye color, textura, bordes, formas y relaciones espaciales. En este tipo de escenarios, el aprendizaje automático tradicional depende de la calidad de la representación construida manualmente, mientras que el aprendizaje profundo ofrece la posibilidad de aprender dichas representaciones directamente desde la imagen \cite{aws2026ml_vs_dl}.

Esta capacidad resulta ser valiosa cuando se trabaja con imágenes recolectadas en condiciones reales, donde existe variabilidad en iluminación, ángulo, fondo y calidad de captura. En dichas condiciones, la posibilidad de que el modelo aprenda representaciones más complejas y adaptativas constituye una ventaja importante frente a enfoques basados exclusivamente en características diseñadas manualmente \cite{aws2026ml_vs_dl}.

En consecuencia, aunque el aprendizaje automático tradicional y el aprendizaje profundo comparten fundamentos relacionados con el aprendizaje a partir de datos, sus diferencias en la forma de representar la información, en el proceso de entrenamiento y en los requerimientos de infraestructura hacen que el aprendizaje profundo resulte más adecuado para problemas visuales complejos. Esta distinción es particularmente relevante para el presente proyecto, donde la entrada del sistema corresponde a imágenes capturadas en campo y el objetivo consiste en identificar patrones visuales asociados a la posible presencia de plaga.

Además, el objetivo del proyecto no se reduce a realizar una clasificación en condiciones ideales, sino a construir una solución que pueda funcionar en un contexto real de uso, orientado a la detección temprana de plagas en árboles. Esto exige un enfoque que sea capaz de responder de forma más consistente ante la diversidad visual presente en las imágenes del caso de estudio. En este sentido, el aprendizaje profundo no solo representa una alternativa técnicamente viable, sino también una elección metodológica coherente con la naturaleza del problema y con las condiciones reales bajo las cuales se espera que opere el sistema.
